\section{Parallel repair branches and cache-aware orchestration}\label{sec:C33:03}
Now allow four isolated workers to investigate different hypotheses. They share a read-only task prefix but have independent writable worktrees. A coordinator receives evidence and candidate patches rather than raw unbounded histories. The meaningful comparison is not simply four agents versus one. Compare against a single agent with the same total budget and four independent samples without communication. Otherwise gains from additional search can be misattributed to coordination.

A precise resource example uses the hypothetical homogeneous model from Chapter~\ref{ch:C31}, with 128 KiB of KV state per retained token. Suppose each worker uses a common 32,768-token prefix followed by a private 8,192-token continuation. Without sharing, four sequences retain 163,840 token positions, consuming 20 GiB of KV state. With exact read-only prefix sharing and private suffix pages, the unique retained positions are $32{,}768+4\times8{,}192=65{,}536$, consuming 8 GiB before overhead. The 12-GiB difference is calculated storage sharing, not a measured speedup or a guarantee that an engine implements it.

Sharing is valid only while the common prefix computation remains compatible. If one worker's context constructor rewrites an early shared observation, its dependent state must diverge. Copy-on-write protects other workers' pages but does not make the edited worker's old suffix valid. A coordinator that changes the common instruction midway through execution creates another compatibility boundary; workers must receive a versioned update or remain explicitly pinned to the old task state.

The wall-clock lower bound is the duration of the longest dependency path, while total resource consumption sums work across branches. A slow branch may contain the only evidence that exposes a constraint violation. Cancelling it solely because another branch produced a plausible patch trades latency against coverage. The selector should know which checks remain incomplete and whether its chosen candidate was examined by an independent validator. A vote among models that all missed the same requirement is not independent verification.

Selection also changes the outcome distribution. Let $E_j$ indicate that worker $j$ has produced an eligible candidate, $T_j$ its completion time, and $V_j$ a true contract violation by that candidate. Conditional on at least one eligible candidate, define $J=\operatorname*{arg\,min}_{j:E_j=1}T_j$ with a fixed tie rule. Then
\begin{equation}
 P(V_J)=\sum_j P(J=j)P(V_j\mid J=j).
\end{equation}
This identity does not allow the conditional risk to be replaced by each worker's marginal risk. Fast shortcuts, early termination, and validation cost can correlate quality with selection. A serving change that changes completion order can therefore change the delivered-artifact distribution even when every worker's policy is fixed. Compare selectors on identical saved candidate sets when isolating selection, and compare full executions when measuring the deployed system. Cancelled branches have unobserved counterfactual completions; counting them as successful or failed would invent outcomes.

The final merge is a separate operation. Two individually passing patches can conflict or invalidate each other's assumptions. Preserve each branch's base commit and diff, merge in an isolated worktree, and rerun integration checks on the merged artifact. If the coordinator adopts only one patch, it should retain the rejected branches' relevant counterevidence. Discarding all unsuccessful trajectories can create a memory bank that remembers only apparent successes and loses the information needed to avoid repeating failed hypotheses.

